% ===== PRÉAMBULE CANONIQUE — à coller VERBATIM en tête de CHAQUE .tex =====
\documentclass[11pt,a4paper]{article}
\usepackage[utf8]{inputenc}\usepackage[T1]{fontenc}\usepackage{lmodern}
\usepackage[french]{babel}
\usepackage{amsmath,amssymb,mathtools}\usepackage{icomma}
\usepackage[version=4]{mhchem}          % \ce{...} : équations & formules
\usepackage{siunitx}\sisetup{locale=FR} % \qty{}{}, \si{}, \ang{}
\usepackage{chemfig}                    % \chemfig{...} : structures développées
\usepackage{xcolor}\usepackage{colortbl}
\usepackage{tikz}\usepackage{pgfplots}\pgfplotsset{compat=1.18}
\usepackage[most]{tcolorbox}\usepackage{enumitem}\usepackage{array,booktabs}
\usepackage[a4paper,margin=2.2cm]{geometry}
\usetikzlibrary{arrows.meta,calc,angles,quotes,positioning,patterns,backgrounds,shapes.geometric}
\definecolor{primary}{RGB}{222,49,99}\definecolor{primarydark}{RGB}{150,22,55}
\definecolor{defcol}{RGB}{0,114,178}\definecolor{methcol}{RGB}{52,92,148}
\definecolor{excol}{RGB}{0,135,90}\definecolor{attncol}{RGB}{200,40,55}
\tcbset{boxbase/.style={enhanced,breakable,boxrule=0.9pt,arc=2.5pt,
  left=3mm,right=3mm,top=2.3mm,bottom=2.3mm,fonttitle=\bfseries,coltitle=white}}
% --- boîtes TUTORIEL (noms EXACTS, ne pas renommer) ---
\newtcolorbox{cle}[1][]{boxbase,colback=defcol!6,colframe=defcol,
  title={\textsc{À retenir}\ifx\relax#1\relax\relax\else\textmd{ : \itshape #1}\fi}}
\newtcolorbox{methode}[1][]{boxbase,colback=methcol!7,colframe=methcol,sharp corners,
  title={\textsc{Méthode}\ifx\relax#1\relax\relax\else\textmd{ --- \itshape #1}\fi}}
\newtcolorbox{exemplebox}[1][]{boxbase,colback=excol!5,colframe=excol,
  title={\textsc{Exemple}\ifx\relax#1\relax\relax\else\textmd{ : \itshape #1}\fi}}
\newtcolorbox{piege}[1][]{boxbase,colback=attncol!6,colframe=attncol,
  title={\textsc{Piège}\ifx\relax#1\relax\relax\else\textmd{ : \itshape #1}\fi}}
\newtcolorbox{remarque}[1][]{boxbase,colback=black!4,colframe=black!45,
  title={\textsc{Remarque}\ifx\relax#1\relax\relax\else\textmd{ : \itshape #1}\fi}}
% --- boîtes EXERCICE / CORRIGÉ (noms EXACTS, auto-comptées) ---
\newtcolorbox[auto counter]{exo}[1][]{boxbase,colback=primary!4,colframe=primary,
  title={\textsc{Exercice}~\thetcbcounter\ifx\relax#1\relax\relax\else\textmd{ --- \itshape #1}\fi}}
\newtcolorbox[auto counter]{corrige}[1][]{boxbase,colback=excol!4,colframe=excol,
  title={\textsc{Corrigé}~\thetcbcounter\ifx\relax#1\relax\relax\else\textmd{ --- \itshape #1}\fi}}
\newcommand{\R}{\mathbb{R}}\newcommand{\N}{\mathbb{N}}\newcommand{\Z}{\mathbb{Z}}\newcommand{\ds}{\displaystyle}
% (un \usepackage{fancyhdr}+en-tête est possible ; il est ignoré côté web)
% =========================================================================
\begin{document}
\sisetup{per-mode=symbol}
\begin{corrige}[très grands, très petits]
On place la virgule après le premier chiffre non nul et l'on compte les rangs.
\begin{enumerate}[label=\alph*)]
  \item $602\,200\,000\,000\,000\,000\,000\,000 = 6,022\times10^{23}$.
  \item $0,0000001 = 1\times10^{-7}\ \si{\mol\per\litre}$ (la virgule avance de $7$ rangs).
  \item $0,000\,000\,000\,154 = 1,54\times10^{-10}\ \si{\metre}$.
\end{enumerate}
\end{corrige}

\begin{corrige}[notation ingénieur et préfixes SI]
On amène l'exposant au multiple de $3$ le plus proche ; il correspond alors à un préfixe.
\begin{enumerate}[label=\alph*)]
  \item $2,096\times10^{-2}\ \si{\litre} = 20,96\times10^{-3}\ \si{\litre} = \qty{20,96}{\milli\litre}$.
  \item $4,7\times10^{-6}\ \si{\metre} = \qty{4,7}{\micro\metre}$.
  \item $1,5\times10^{3}\ \si{\gram} = \qty{1,5}{\kilo\gram}$.
\end{enumerate}
\end{corrige}

\begin{corrige}[arrondir avec retenue]
On garde $n$ CS ; l'arrondi par excès d'un $9$ propage une retenue.
\begin{enumerate}[label=\alph*)]
  \item $0,0995 \to 0,10$ (soit $1,0\times10^{-1}$) : chiffre supprimé $5\ge5$, la retenue remonte.
  \item $7,996 \to 8,00$.
  \item $149,7 \to 150$, écrit $1,50\times10^{2}$ pour afficher $3$ CS.
  \item $3,0499 \to 3,05$.
\end{enumerate}
\end{corrige}

\begin{corrige}[n'arrondir qu'une seule fois]
\begin{enumerate}[label=\alph*)]
  \item Il faut arrondir \textbf{une seule fois}, directement à partir du nombre complet : on regarde le
        premier chiffre supprimé du nombre \emph{de départ}, pas d'un résultat déjà arrondi.
  \item $2,748$ à $2$ CS : le premier chiffre supprimé est $4<5$, donc $\boxed{2,7}$. La méthode en deux
        temps donne $2,8$, ce qui est \emph{faux} (elle a compté deux fois une retenue inexistante).
\end{enumerate}
\end{corrige}

\begin{corrige}[lire l'ordre de grandeur]
On compare d'abord les puissances de dix, puis les mantisses à exposant égal :
\[ 8,1\times10^{-4}\ <\ 1,2\times10^{-3}\ <\ 3,4\times10^{-3}\ <\ 2,9\times10^{-2}. \]
\end{corrige}

\end{document}
